\documentclass{IOS-Book-Article}

\usepackage{times}
\usepackage{soul}
\usepackage{graphicx}
\usepackage{amsmath,amssymb,amsfonts}
\usepackage{booktabs}
\usepackage{multirow}
\usepackage{tabularx}
\usepackage{microtype}

\newtheorem{definition}{Definition}
\newtheorem{example}{Example}
\newtheorem{theorem}{Theorem}

% -------------------------------------
% Document Start
% -------------------------------------
\begin{document}

\begin{frontmatter}

\title{Measuring the Downstream Decision Impact of Predicate Extraction Errors in Symbolic Legal Engines}
\runningtitle{Measuring the Downstream Decision Impact of Predicate Extraction Errors in Symbolic Legal Engines}

% Authors & Affiliations
\author[A]{\fnms{Linh} \snm{Pham Van}%
\thanks{Corresponding Author: Linh Pham Van, College of Engineering \& Computer Science, VinUniversity, Viet Nam; E-mail: linhpv.vinuni@gmail.com.}},
\author[A]{\fnms{Ha-Thanh} \snm{Nguyen}%
\thanks{Corresponding Author: Ha-Thanh Nguyen, College of Engineering \& Computer Science, VinUniversity, Viet Nam; E-mail: thanh.nh6@vinuni.edu.vn.}}
and
\author[B]{\fnms{Ken} \snm{Satoh}%
\thanks{Corresponding Author: Ken Satoh, Center for Juris-Informatics, ROIS-DS, Japan; E-mail: ksatoh@nii.ac.jp.}}
\runningauthor{L.V. Pham et al.}

\address[A]{College of Engineering \& Computer Science, VinUniversity, Viet Nam}
\address[B]{Center for Juris-Informatics, ROIS-DS, Japan}

% Abstract
\begin{abstract}
Neuro-symbolic legal reasoning pipelines couple an LLM, which extracts structured predicates from a fact pattern, with a symbolic rule engine that reaches an auditable decision. Such pipelines are typically evaluated by predicate-level F1, which treats every predicate error as equally consequential - an assumption that fails once a rule tree's conditional structure lets some errors flip the decision while others are absorbed as redundant. Using the PYTHEN framework on LegalBench's hearsay classification task, we introduce a predicate-level intervention procedure that corrects one erroneous predicate at a time, reruns the rule tree, and measures whether the decision changes, yielding a Decision Impact Rate (DIR) per predicate, which we compare against standalone extraction F1. Across three LLMs, extraction difficulty and decision impact diverge sharply - the hardest predicate to extract is also the most decision-critical, while a comparably hard predicate is almost always inconsequential, and nominally identical error types carry markedly different downstream risk across models - arguing for rule-aware error attribution as a necessary complement to aggregate predicate accuracy when evaluating neuro-symbolic legal reasoning pipelines.
\end{abstract}

% Keywords
\begin{keyword}
Neuro-Symbolic Legal AI \sep Explainability and Verification \sep Error Attribution \sep Large Language Models \sep LegalBench \sep Hearsay Classification
\end{keyword}

\end{frontmatter}

% -------------------------------------
% Sections
% -------------------------------------
\section{Introduction}
\label{sec:intro}

Neuro-symbolic legal reasoning pipelines couple an LLM, which extracts structured predicates from a fact pattern, with a symbolic rule engine that applies an auditable rule to reach a decision~\cite{yang2025neuro}. Such systems are attractive because the rule structure is fixed and inspectable, in principle allowing the decision path to be verified independently of the model that produced it. We instantiate this division of labour as PYTHEN, a neuro-symbolic pipeline for hearsay classification under LegalBench~\cite{guha2023legalbench}, and show that it improves F1-score and reduces instability across model families relative to both direct and chain-of-thought prompting. An oracle experiment, in which gold predicates are fed directly into the rule tree, further shows that the symbolic component reaches perfect agreement with ground truth once its inputs are correct - indicating that essentially all residual error of the pipeline originates upstream, in the LLM's predicate extraction step.

This result exposes a gap that standard evaluation practice does not close. Predicate extraction is typically assessed by aggregate precision/recall/F1~\cite{olek2023evaluationf1scorerecent}, which implicitly treats every predicate error as equally consequential. This assumption fails once predicates are consumed by a rule tree: a rule tree encodes conditional structure, so some predicates sit on a decisive branch where a single wrong value flips the final verdict, while others are redundant with correlated conditions, short-circuited by an earlier branch, or on a path the specific instance never reaches. Two models with identical predicate-level F1 can therefore differ substantially in how many of their errors actually reach the final decision, and aggregate extraction metrics cannot distinguish the two cases, because they are computed independently of the rule structure the predicates feed into.

This paper asks, for each predicate error an LLM extractor makes, not only \emph{that} it occurred but \emph{whether it matters} - whether correcting it would change the system's output. We introduce a predicate-level intervention procedure: for each fact pattern, we correct one erroneous predicate at a time toward its gold value, rerun the PYTHEN rule tree, and observe whether the decision changes. This yields a Decision Impact Rate (DIR) for every predicate, which we compare against standalone extraction F1, and use to refine an extraction-error taxonomy along a decision-relevance dimension. Across three LLMs, we find that extraction difficulty and decision impact diverge sharply - the hardest predicate to extract is also the most decision-critical, while a comparably hard predicate is almost always inconsequential - and that nominally identical error types carry markedly different downstream risk across models. These results argue for rule-aware error attribution, rather than aggregate predicate accuracy alone, as a necessary complement when evaluating and verifying neuro-symbolic legal reasoning pipelines.

\section{Related Work}
\label{sec:related}

Formalizing statutory rules as executable logic dates to Sergot et al.'s encoding of the British Nationality Act as a Prolog program~\cite{sergot1986}, and continues through systems such as PROLEG~\cite{satoh2010proleg} and defeasible-argumentation frameworks~\cite{benchcapon2007,prakken2020}, whose soundness and explainability depend on a correctly and completely specified rule structure - traditionally hand-entered rather than extracted from unstructured text. The rise of LLMs has shifted attention toward automatically bridging this gap. LegalBench~\cite{guha2023legalbench} evaluates LLMs on legal reasoning tasks via direct or chain-of-thought prompting for a final label, an evaluation paradigm whose reliability is questionable: chain-of-thought explanations can be systematically unfaithful to the process that actually produced a model's answer~\cite{turpin2023unfaithful}, and a related faithfulness study of LLMs as autoformalizers in legal reasoning documents an analogous divergence between benchmark accuracy and logical faithfulness~\cite{wang2026knowyourlimits}. Closer to this paper's architecture, several systems couple LLMs with an explicit symbolic component rather than relying on the LLM's final-answer judgment alone, e.g., knowledge-graph-guided legal provision prediction~\cite{li2021legalkg} and a multi-agent framework that formalizes statutes into verifiable intermediate representations before symbolic inference~\cite{sadowski2025verifiable}. These works establish that routing LLM output through a symbolic intermediate is broadly beneficial, consistent with what this paper also finds, but they evaluate only whether the final decision is correct, not which intermediate extraction errors were responsible when it is not - the question this paper addresses through predicate-level intervention on a fixed, auditable rule tree.

\section{Methodology}
\label{sec:method}

We formalize the neuro-symbolic pipeline coupling an LLM predicate extractor with the PYTHEN rule tree, and the predicate-level intervention procedure used to measure the downstream decision impact of individual predicate extraction errors.

\textbf{Task, pipeline, and rule tree.}
We use the \texttt{hearsay} subset of LegalBench (94 fact patterns), decomposed under Federal Rule of Evidence 801 into five binary predicates - \textit{is\_oral\_assertion}, \textit{is\_written\_assertion}, \textit{is\_non\_verbal\_assertive\_conduct}, \textit{is\_out\_of\_court}, and \textit{is\_introduced\_to\_prove\_truth} - which we manually annotated for every instance, giving a gold table of $94\times5=470$ predicate slots. Given a fact pattern $x$, an LLM extractor $\mathcal{E}$ produces $\hat{p} = \mathcal{E}(x) \in \{0,1\}^5$, which is passed unchanged to a fixed, hand-authored rule tree $R$ encoding FRE~801:
\begin{equation}
\label{eq:hearsay-rule}
R(p) = \big(p_{\text{oral}} \lor p_{\text{written}} \lor p_{\text{nonverbal}}\big) \land\, p_{\text{out-of-court}} \land\, p_{\text{truth}},
\end{equation}
yielding $\hat{y} = R(\hat{p}) \in \{\text{Hearsay}, \text{Not Hearsay}\}$. Because $R$ is fixed and independent of the LLM, any difference in output across models is fully attributable to $\mathcal{E}$. The disjunctive clause and the two conjunctive gates in Eq.~\ref{eq:hearsay-rule} make predicates unequally important for a given instance: predicate $p_j$ is \emph{structurally critical} if flipping it alone, with all other predicates held at their gold values, changes $R(p)$, and \emph{structurally redundant} otherwise - e.g., $p_{\text{written}}$ is redundant whenever $p_{\text{oral}}$ or $p_{\text{nonverbal}}$ already holds, whereas $p_{\text{out-of-court}}$ and $p_{\text{truth}}$ are always critical when the assertion clause holds, since either alone can veto the decision.

\textbf{Predicate-level intervention and decision impact.}
Standard evaluation compares $\hat{y}_i = R(\hat{p}_i)$ against gold $y_i$ and separately reports predicate accuracy, but neither reveals which of the (possibly several) simultaneous predicate errors on an instance caused a wrong decision. For each instance $i$ with erroneous predicate indices $E_i = \{j : \hat{p}_{i,j} \neq p_{i,j}\}$, and each $j \in E_i$, we correct only predicate $j$ to its gold value, leaving every other predicate - including any other errors - at its extracted value, and rerun the unmodified rule tree to obtain $\hat{y}_i^{(j)}$. This isolates the causal effect of correcting exactly one predicate error, requiring no retraining or re-prompting: only the symbolic stage is rerun. The \emph{decision impact} of that single error is the binary indicator
\begin{equation}
\label{eq:impact-indicator}
\mathrm{Impact}(i,j) = \mathbb{1}\big[\hat{y}_i^{(j)} \neq \hat{y}_i\big], \qquad j \in E_i,
\end{equation}
and averaging Eq.~\ref{eq:impact-indicator} over all observed errors of predicate $j$ across the dataset gives its \emph{Decision Impact Rate},
\begin{equation}
\label{eq:dir}
\mathrm{DIR}(j) = \frac{1}{|\{i : j \in E_i\}|} \sum_{i \,:\, j \in E_i} \mathrm{Impact}(i,j),
\end{equation}
which we compare against predicate $j$'s standalone extraction $F_1$ to test whether extraction difficulty predicts practical importance to the downstream decision. We further annotate every extraction error with a type - E1: Predicate Hallucination, E2: Predicate Omission, E3: Invalid Reference - assigned from the extracted JSON, the gold table, and a supporting text span the extractor is prompted to emit for each true predicate, and with the impact label (Critical/Benign) from Eq.~\ref{eq:impact-indicator}, yielding six refined error categories used in Section~\ref{sec:experiments} to determine, per model, what fraction of extraction errors materially undermine the reliability of the final decision.

\section{Experiments and Evaluation}
\label{sec:experiments}

We evaluate three LLMs as predicate extractors - DeepSeek V3.2, GPT-OSS 120B, and Gemma 4 31B - on the 94-instance \texttt{hearsay} subset, under three configurations: \textbf{Naive} direct prompting for a binary verdict, \textbf{Chain-of-Thought (CoT)}, and \textbf{Neuro-symbolic (Ours)}, where the LLM only extracts the five predicates for the PYTHEN rule tree (Eq.~\ref{eq:hearsay-rule}). The Neuro-symbolic pipeline improves F1-score over both baselines for all three models (e.g., from 0.725 under CoT to 0.813 for GPT-OSS 120B) and is the only configuration in which every model's accuracy exceeds 0.77. Feeding the manually annotated gold predicate table directly into the rule tree, bypassing $\mathcal{E}$ entirely, yields Accuracy = Precision = Recall = $F_1 = 1.0$: the rule tree is logically sound with respect to this dataset, so all residual error of the Neuro-symbolic pipeline is attributable to the extraction step $\mathcal{E}$, licensing the intervention analysis below.

Table~\ref{tab:dir} applies the intervention procedure to every predicate error produced by each model, reporting the standalone extraction $F_1$ (over all $470$ predicate slots) alongside the resulting Decision Impact Rate (DIR); subscripts give the number of erroneous slots each DIR is computed over. \textit{is\_written\_assertion} is the hardest predicate to extract for every model ($F_1 \in [0.286, 0.429]$) and simultaneously has, by a wide margin, the highest DIR (0.700-0.900): nearly every time a model gets this predicate wrong, that error flips the final decision. \textit{is\_non\_verbal\_assertive\_conduct} has comparably low $F_1$ (0.593-0.667) but the lowest DIR of all five predicates (0.000-0.091), since its errors are almost always structurally redundant, absorbed whenever \textit{is\_oral\_assertion} or \textit{is\_written\_assertion} already satisfies the disjunctive clause. The two conjunctive gates, \textit{is\_out\_of\_court} and \textit{is\_introduced\_to\_prove\_truth}, are extracted comparatively well yet carry disproportionate DIR (up to 1.000 and 0.875, respectively) because a conjunctive condition has no disjunctive backup to absorb a mistake. Overall, extraction difficulty ($1-F_1$) correlates weakly or negatively with DIR for four of the five predicates, and strongly positively only for \textit{is\_written\_assertion} - a predicate that is both hard to extract and disproportionately consequential when wrong, making it the clearest target for extractor improvement.

\begin{table}[htbp]
\centering
\small
\caption{Decision Impact Rate (DIR) versus extraction F1 by predicate and model. Subscripts denote the number of observed errors the DIR is computed over.}
\label{tab:dir}
\begin{tabular}{lcccccc}
\toprule
& \multicolumn{2}{c}{\textbf{DeepSeek V3.2}} & \multicolumn{2}{c}{\textbf{GPT-OSS 120B}} & \multicolumn{2}{c}{\textbf{Gemma 4 31B}} \\
\textbf{Predicate} & F1 & DIR & F1 & DIR & F1 & DIR \\
\midrule
is\_oral\_assertion                 & 0.897 & $0.182_{(11)}$ & 0.929 & $0.125_{(8)}$  & 0.957 & $0.200_{(5)}$  \\
is\_written\_assertion              & 0.286 & $\mathbf{0.700_{(10)}}$ & 0.429 & $\mathbf{0.875_{(8)}}$  & 0.375 & $\mathbf{0.900_{(10)}}$ \\
is\_non\_verbal\_assertive\_conduct & 0.593 & $0.091_{(11)}$ & 0.667 & $0.000_{(9)}$  & 0.667 & $0.000_{(8)}$  \\
is\_out\_of\_court                  & 0.910 & $0.167_{(12)}$ & 0.877 & $0.188_{(16)}$ & 0.976 & $1.000_{(3)}$  \\
is\_introduced\_to\_prove\_truth    & 0.760 & $0.621_{(29)}$ & 0.744 & $0.394_{(33)}$ & 0.864 & $0.875_{(16)}$ \\
\bottomrule
\end{tabular}
\end{table}

Table~\ref{tab:taxonomy} applies the impact label to every E1-E3 extraction error over the 470 predicate slots, showing that the same nominal error type carries different practical risk across models. Hallucinated predicates (E1) are decision-critical only 59.5\% and 46.9\% of the time for DeepSeek V3.2 and GPT-OSS 120B, respectively - close to half of these errors are absorbed by the rule tree's disjunctive structure and never affect the output - whereas 96.2\% of Gemma 4 31B's E1 errors are decision-critical, despite hallucinating about as often in absolute terms, because it rarely hallucinates an already-redundant predicate. Omission errors (E2) are critical far less often across all three models (4.0-13.8\%), indicating that a missed predicate is disproportionately likely to be one the decision did not depend on - an asymmetry not visible from raw E1/E2 frequency counts alone. E3 (invalid reference) is negligible in frequency (at most 2 slots) and is included in the table for completeness only.

\begin{table}[htbp]
\centering
\small
\caption{Refined error taxonomy: critical vs.\ benign extraction errors by model, out of 470 predicate slots.}
\label{tab:taxonomy}
\begin{tabular}{llccc}
\toprule
\textbf{Error Type} & \textbf{Impact} & \textbf{DeepSeek V3.2} & \textbf{GPT-OSS 120B} & \textbf{Gemma 4 31B} \\
\midrule
\multirow{3}{*}{E1: Hallucination} & Critical & 25 (5.32\%) & 23 (4.89\%) & 25 (5.32\%) \\
                                    & Benign   & 17 (3.62\%) & 26 (5.53\%) & 1 (0.21\%) \\
                                    & \% Critical & 59.5\% & 46.9\% & \textbf{96.2\%} \\
\midrule
\multirow{3}{*}{E2: Omission}      & Critical & 4 (0.85\%) & 1 (0.21\%) & 2 (0.43\%) \\
                                    & Benign   & 25 (5.32\%) & 24 (5.11\%) & 14 (2.98\%) \\
                                    & \% Critical & 13.8\% & 4.0\% & 12.5\% \\
\midrule
\multirow{3}{*}{E3: Invalid Ref.}   & Critical & 1 (0.21\%) & 0 (0.00\%) & 0 (0.00\%) \\
                                    & Benign   & 1 (0.21\%) & 0 (0.00\%) & 0 (0.00\%) \\
                                    & \% Critical & 50.0\% & - & - \\
\bottomrule
\end{tabular}
\end{table}

\section{Discussion}
\label{sec:discussion}

Predicate extraction errors do not affect the final decision equally, and their impact depends on where a predicate sits in the rule structure - invisible to aggregate F1 or accuracy. This supports rule-aware verification: flag only predicates that are both error-prone and high-DIR for the instance's rule path - \textit{is\_written\_assertion} and \textit{is\_introduced\_to\_prove\_truth} in the hearsay setting, and \textit{is\_out\_of\_court} when the assertion clause is otherwise unsatisfied - while deprioritizing low-DIR predicates such as \textit{is\_non\_verbal\_assertive\_conduct}. The taxonomy results add that the same error type carries different risk across models: Gemma 4 31B's hallucinated predicates are decision-critical 96.2\% of the time, versus 46.9-59.5\% for the other two, despite comparable aggregate accuracy - a gap a model-selection decision based on F1 alone would miss.

\textbf{Limitations.} All results come from a single LegalBench task with a shallow rule tree; whether this divergence generalizes to deeper or exception-laden rule structures, and to joint (multi-predicate) intervention, is left to future work. Some DIR estimates rest on as few as 3-5 errors and are directionally informative but statistically fragile, and both the gold table and the rule tree's soundness are validated only within this dataset.

\section{Conclusion}
\label{sec:conclusion}

Correcting one erroneous predicate at a time and rerunning the PYTHEN rule tree, we introduced a Decision Impact Rate (DIR) that diverges sharply from standalone extraction F1 across three LLMs: \textit{is\_written\_assertion} is both hardest to extract and most decision-critical, \textit{is\_non\_verbal\_assertive\_conduct} is hard to extract yet redundant, and identical error types carry different risk across models. This supports rule-aware error attribution, rather than aggregate accuracy alone, for evaluating neuro-symbolic legal pipelines.


% -------------------------------------
% References
% -------------------------------------
\bibliographystyle{vancouver}
\bibliography{references}

\end{document}
